% Primer teorema de isomorfía

\begin{frame}
    \frametitle{Un ejemplo antes del teorema}

    \begin{columns}[c]
        \begin{column}{0.55\textwidth}
            \begin{example}
                Sea $f\colon\R^2\to\R$, $f(u_1,u_2)=u_1-u_2$. Es sobreyectiva
                y $\ker f=\langle(1,1)\rangle$.
            \end{example}

            Para cada $c\in\R$, la \alert{fibra}
            \[
                f^{-1}(c)=\{u:\ u_1-u_2=c\}
            \]
            es una recta paralela a $\ker f$, es decir, una clase módulo
            $\ker f$.

            \medskip

            Cada número real $c$ corresponde a una única clase de
            $\R^2/\ker f$, y viceversa: $f$ identifica el cociente con
            $\R=\im f$.
        \end{column}
        \begin{column}{0.42\textwidth}
            \centering
            \begin{tikzpicture}[scale=0.7]
                \clip (-3.2,-2.6) rectangle (3.4,2.8);
                \draw[->] (-2.5,0) -- (2.5,0) node[below left] {\footnotesize $u_1$};
                \draw[->] (0,-2.5) -- (0,2.5) node[below left] {\footnotesize $u_2$};
                \foreach \c in {-2,-1,1,2}
                    \draw[gray!60] (-3,-3-\c) -- (3,3-\c);
                \draw[very thick,alertcolor] (-3,-3) -- (3,3);
                \foreach \c in {0,1,2}
                    \node[right,font=\footnotesize,fill=white,inner sep=1pt] at (2.5,2.5-\c) {$\c$};
                \foreach \c in {-2,-1}
                    \node[left,font=\footnotesize,fill=white,inner sep=1pt] at (-2.5,-2.5-\c) {$\c$};
            \end{tikzpicture}
        \end{column}
    \end{columns}
\end{frame}


\begin{frame}
    \frametitle{Las fibras son las clases módulo el núcleo}

    \begin{lemma}
        Sea $f\colon V\to V'$ un homomorfismo. Para todo $u,u'\in V$,
        \[
            f(u)=f(u')\iff u+\ker f=u'+\ker f.
        \]
    \end{lemma}

    \begin{proof}
        Por la caracterización de la inyectividad,
        $f(u)=f(u')\iff u-u'\in\ker f$, y por la caracterización de las
        clases módulo un subespacio (Tema~3),
        $u-u'\in\ker f\iff u+\ker f=u'+\ker f$.
    \end{proof}

    \medskip

    Es decir, $f$ \alert{no distingue} los vectores de una misma clase
    módulo $\ker f$, y distingue los de clases distintas.
\end{frame}


\begin{frame}[allowframebreaks, t]
    \frametitle{Primer teorema de isomorfía}

    \begin{theorem}
        Sea $f\colon V\to V'$ un homomorfismo. La aplicación
        \[
            \bar f\colon V/\ker f\longrightarrow\im f,
            \qquad
            \bar f(u+\ker f)=f(u),
        \]
        está bien definida y es un isomorfismo. En particular,
        \[
            V/\ker f\cong\im f.
        \]
    \end{theorem}

    \framebreak

    \begin{proof}
        Escribamos $K=\ker f$. Por el lema, para todo $u,u'\in V$,
        \[
            u+K=u'+K\iff f(u)=f(u').
        \]
        La implicación ($\Rightarrow$) dice que $\bar f$ \alert{está bien
        definida}: el valor $f(u)$ no depende del representante $u$ de la
        clase. La implicación ($\Leftarrow$) dice que $\bar f$ es
        \alert{inyectiva}.

        \alert{Lineal.} Para $u,u'\in V$ y $\lambda\in k$,
        \[
            \bar f\bigl((u+K)+(u'+K)\bigr)=\bar f\bigl((u+u')+K\bigr)
            =f(u+u')=\bar f(u+K)+\bar f(u'+K),
        \]
        \[
            \bar f\bigl(\lambda(u+K)\bigr)=\bar f(\lambda u+K)
            =f(\lambda u)=\lambda\bar f(u+K).
        \]

        \alert{Sobreyectiva.} Todo elemento de $\im f$ es $f(u)=\bar f(u+K)$
        para algún $u\in V$.
    \end{proof}
\end{frame}


\begin{frame}[fragile]
    \frametitle{Factorización canónica de un homomorfismo}

    Todo homomorfismo $f\colon V\to V'$ se descompone como
    \[
        \begin{tikzcd}[column sep=large]
            V\arrow[rr,"f"]\arrow[dr,"\pi"'] & & V'\\
            & V/\ker f\arrow[r,"\bar f"',"\cong"] & \im f\arrow[u,hook,"\iota"']
        \end{tikzcd}
    \]
    es decir, $f=\iota\circ\bar f\circ\pi$, donde
    \begin{itemize}
        \item $\pi$ es la proyección canónica, \alert{sobreyectiva};
        \item $\bar f$ es un \alert{isomorfismo};
        \item $\iota$ es la inclusión, \alert{inyectiva}.
    \end{itemize}

    \medskip

    En efecto, $\iota\bar f\pi(u)=\iota\bar f(u+\ker f)=f(u)$.

    \medskip

    Así, \alert{todo} homomorfismo es, salvo isomorfismo, una proyección
    sobre un cociente seguida de una inclusión.
\end{frame}


\begin{frame}
    \frametitle{Consecuencias}

    \begin{corollary}
        Sea $f\colon V\to V'$ un homomorfismo.
        \begin{enumerate}[(1)]
            \item Si $f$ es sobreyectivo, $V'\cong V/\ker f$.
            \item Si $V$ tiene dimensión finita, $\dim V=\dim\ker f+\dim\im f$.
        \end{enumerate}
    \end{corollary}

    \begin{proof}
        (1) Es el teorema con $\im f=V'$.
        (2) Como $V/\ker f\cong\im f$ y los espacios isomorfos tienen la
        misma dimensión, $\dim\im f=\dim V/\ker f=\dim V-\dim\ker f$
        (Tema~3).
    \end{proof}

    \medskip

    Es una segunda demostración del teorema de la dimensión.

    \medskip

    \alert{Los cocientes de $V$ son, salvo isomorfismo, las imágenes de
    los homomorfismos que salen de $V$}: un espacio $V'$ es isomorfo a un
    cociente de $V$ si y solo si existe un homomorfismo sobreyectivo
    $V\to V'$.
\end{frame}


\begin{frame}
    \frametitle{Identificar cocientes: ejemplos del Tema 3}

    El teorema explica los cocientes que calculamos en el Tema~3.

    \begin{example}
        \begin{enumerate}
            \item $\R^2/\langle(1,1)\rangle\cong\R$, mediante
            $f(u_1,u_2)=u_1-u_2$: la clase de $(u_1,u_2)$ queda
            determinada por $u_1-u_2$.
            \item Si $W=\{p\in\R_2[x]:\ p(-1)=0\}$, entonces
            $\R_2[x]/W\cong\R$, mediante $p\mapsto p(-1)$: es un
            homomorfismo sobreyectivo cuyo núcleo es $W$.
            \item Si $W=\langle(1,0,1,0),(0,1,0,1)\rangle\subseteq\R^4$,
            entonces $\R^4/W\cong\R^2$, mediante
            $f(u_1,u_2,u_3,u_4)=(u_1-u_3,\ u_2-u_4)$, que es sobreyectivo
            con $\ker f=W$.
            \item $k[x]/\{p:\ p(0)=0\}\cong k$, mediante $p\mapsto p(0)$,
            y $\dim=1$ aunque $k[x]$ tenga dimensión infinita.
            \item La traza $\operatorname{tr}\colon M_n(k)\to k$ es
            sobreyectiva, con núcleo las matrices de traza nula:
            $M_n(k)/\{X:\ \operatorname{tr}X=0\}\cong k$.
        \end{enumerate}
    \end{example}

    En todos los casos, $\bar f$ es la aplicación «calcular el valor
    que recuerda el cociente».
\end{frame}


\begin{frame}
    \frametitle{Ejemplo: un cociente de matrices}

    \begin{example}
        Sea $S\subseteq M_2(\R)$ el subespacio de las matrices simétricas
        y $f\colon M_2(\R)\to M_2(\R)$, $f(X)=X-X^t$.
    \end{example}

    Como
    \[
        f\begin{pmatrix}a&b\\c&d\end{pmatrix}
        =\begin{pmatrix}0&b-c\\c-b&0\end{pmatrix},
    \]
    se tiene $\ker f=S$ e $\im f=\langle E_{12}-E_{21}\rangle$. Por
    el primer teorema de isomorfía,
    \[
        M_2(\R)/S\cong\langle E_{12}-E_{21}\rangle\cong\R,
    \]
    y la clase de $X$ queda determinada por $b-c$.
\end{frame}


\begin{frame}
    \frametitle{Un cociente es isomorfo a cualquier suplementario}

    \begin{proposition}
        Sean $W$ y $S$ subespacios de $V$ con $V=W\oplus S$. La restricción
        a $S$ de la proyección canónica,
        \[
            \pi|_S\colon S\longrightarrow V/W,\qquad s\mapsto s+W,
        \]
        es un isomorfismo. En particular, $S\cong V/W$.
    \end{proposition}

    \begin{proof}
        $\ker(\pi|_S)=S\cap\ker\pi=S\cap W=\{0\}$, luego $\pi|_S$ es
        inyectiva. Es sobreyectiva porque cada $u\in V$ es $u=w+s$ con
        $w\in W$ y $s\in S$, y entonces
        $u+W=s+W=\pi|_S(s)$.
    \end{proof}

    \medskip

    En particular, \alert{todos los suplementarios de $W$ son isomorfos
    entre sí}, y a $V/W$. Por ejemplo, para $W=\langle(1,0,1,0),(0,1,0,1)\rangle$
    en $\R^4$, tanto $\langle e_1,e_2\rangle$ como $\langle e_3,e_4\rangle$
    son suplementarios y son isomorfos a $\R^4/W\cong\R^2$.
\end{frame}


\begin{frame}
    \frametitle{Factorización a través de un cociente}

    \begin{proposition}
        Sea $f\colon V\to V'$ un homomorfismo y sea $W$ un subespacio de
        $V$ con $W\subseteq\ker f$. La aplicación
        \[
            \bar f\colon V/W\longrightarrow V',\qquad
            \bar f(u+W)=f(u),
        \]
        está bien definida y es un homomorfismo, con
        $\bar f\circ\pi=f$, $\im\bar f=\im f$ y
        $\ker\bar f=\{u+W:\ u\in\ker f\}$.
    \end{proposition}

    \begin{proof}
        Si $u+W=u'+W$, entonces $u-u'\in W\subseteq\ker f$ y
        $f(u)=f(u')$: está bien definida. La linealidad se comprueba como
        en el primer teorema de isomorfía. Las demás afirmaciones son
        inmediatas: $\bar f(u+W)=0$ si y solo si $f(u)=0$.
    \end{proof}

    \medskip

    Un homomorfismo definido en $V$ que se anula en $W$ \alert{pasa al
    cociente} $V/W$. El primer teorema de isomorfía es el caso
    $W=\ker f$.
\end{frame}


\begin{frame}
    \frametitle{Segundo teorema de isomorfía}

    \begin{theorem}
        Sean $U$ y $W$ subespacios de $V$. Entonces
        \[
            U/(U\cap W)\cong(U+W)/W.
        \]
    \end{theorem}

    \begin{proof}
        Sea $\varphi\colon U\to(U+W)/W$, $\varphi(u)=u+W$. Es un
        homomorfismo (composición de la inclusión $U\subseteq U+W$ y la
        proyección canónica). Es sobreyectivo: toda clase de $(U+W)/W$ es
        $(u+w)+W=u+W=\varphi(u)$ con $u\in U$, $w\in W$. Y
        \[
            \ker\varphi=\{u\in U:\ u+W=0\}=U\cap W.
        \]
        Por el primer teorema de isomorfía, $U/(U\cap W)\cong(U+W)/W$.
    \end{proof}
\end{frame}


\begin{frame}[allowframebreaks, t]
    \frametitle{La fórmula de Grassmann, de nuevo}

    \begin{corollary}[Fórmula de Grassmann]
        Si $U$ y $W$ son subespacios de dimensión finita de $V$,
        \[
            \dim(U+W)=\dim U+\dim W-\dim(U\cap W).
        \]
    \end{corollary}

    \begin{proof}
        Por el segundo teorema, los dos espacios tienen la misma
        dimensión, y por la dimensión del cociente (Tema~3),
        \[
            \dim U-\dim(U\cap W)=\dim(U+W)-\dim W.\qedhere
        \]
    \end{proof}

    \framebreak

    \begin{example}
        Si $U$ y $W$ son dos planos distintos de $\R^3$, entonces
        $U+W=\R^3$ y $U\cap W$ es una recta, y el teorema da
        \[
            U/(U\cap W)\cong\R^3/W\cong\R,
        \]
        ambos de dimensión $2-1=3-2=1$.
    \end{example}
\end{frame}
